\chapter{Options à zéro jour, hebdomadaires et flux d'options de détail}\label{ch:m1:zero-day-options-and-retail-flow}

À deux heures de l'après-midi, l'option la plus activement échangée sur l'indice
boursier le plus important du monde n'a plus que deux heures à vivre. Elle a été cotée il y a
quelques semaines, peu échangée jusqu'à ce matin, et à la clôture elle vaudra
soit zéro, soit la distance entre l'indice et son prix d'exercice. En
2025, les options arrivant à échéance le jour même ont représenté en moyenne 2.3 millions de contrats par jour
sur les options sur l'indice S\&P 500, soit 59~\% de tout leur volume. Une option aussi proche de
l'échéance est un instrument différent du contrat à trois mois décrit au
\cref{ch:m1:volatility-first-contact}~: presque insensible à la \omterm{def:m1:volatility-first-contact:iv}{volatilité implicite},
extrêmement sensible à la trajectoire du sous-jacent dans les minutes
qui suivent, et couverte, par la contrepartie, par une quantité
d'actions qui peut passer de rien à tout pour un mouvement
d'un demi pour cent. Ce chapitre mesure cette variation.

\section{Des mensuelles aux quotidiennes}

\begin{definition}[Options hebdomadaires et options à zéro jour]\label{def:m1:zero-day-options-and-retail-flow:zero}
Une \emph{option hebdomadaire}\index{option hebdomadaire} est une \omterm{def:m1:option-contracts-and-exchanges:series}{série d'options} cotées
dont l'échéance n'est pas celle, standard, du mois, et qui est cotée quelques semaines avant
d'arriver à échéance. Une \emph{option à zéro jour}\index{option à zéro jour} (0DTE, zero days to
expiry) est toute option le jour même de son échéance~: ce n'est pas un produit distinct, mais
le dernier jour de vie d'une série hebdomadaire, mensuelle ou trimestrielle.
\end{definition}

\begin{dated}{2026-09}{Comment une échéance chaque jour est apparue}\label{dat:m1:zero-day-options-and-retail-flow:dailies}
Les \omterm{def:m1:zero-day-options-and-retail-flow:zero}{options hebdomadaires} sur l'indice S\&P 500 n'arrivaient d'abord à échéance que le vendredi~;
les échéances du lundi et du mercredi sont venues ensuite. En 2022, la bourse a ajouté les échéances du mardi, à partir du 18 avril, et
celles du jeudi, à partir du 11 mai~: depuis lors, une série arrive à échéance
chaque jour de séance. En 2025, les options à échéance le jour même ont représenté en moyenne 2.3 millions de contrats par jour,
soit 59~\% du volume des options sur indices, sur un marché américain des options qui a échangé
15.2 milliards de contrats dans l'année, soit 26~\% de plus qu'en 2024.
\end{dated}

L'attrait tient à l'arithmétique. Une option \omterm{def:m1:volatility-first-contact:atm}{à la monnaie} avec un jour à courir coûte
environ $0.8\,\sigma\sqrt{T}$ du sous-jacent, soit quelque 0.8~\% pour un indice à 16~\% de
volatilité~: une petite prime pour une position qui double ou disparaît à la
clôture. Les acheteurs s'en servent pour un événement (une publication de données à 08 h 30
ou une décision de banque centrale à 14 h 00) ou comme stop intrajournalier bon marché~; les vendeurs encaissent une prime qui
se déprécie jusqu'à zéro en quelques heures, de manière systématique et en quantité importante.
La \Cref{fig:m1:zero-day-options-and-retail-flow:straddle} montre le point de vue de l'acheteur.

\begin{omfigure}
\begin{tikzpicture}
\begin{axis}[width=10cm,height=4.6cm,
  xlabel={mouvement de l'indice à la clôture (\%)},ylabel={P\&L (\% de l'indice)},
  tick label style={font=\small},label style={font=\small},
  xmin=-2.5,xmax=2.5,ymin=-1.4,ymax=1.9,grid=major,grid style={gray!25},table/col sep=comma]
\addplot[omDef,very thick] table[x=move_pct,y=pnl_pct]{figdata/markets-1/26-zero-day-options-and-retail-flow/straddle.csv};
\draw[thin] (axis cs:-2.5,0) -- (axis cs:2.5,0);
\node[font=\footnotesize,anchor=north] at (axis cs:0,-0.88) {prime 0.80~\%};
\end{axis}
\end{tikzpicture}

\omcaption{Achat du call et du put \omterm{def:m1:volatility-first-contact:atm}{à la monnaie} qui arrivent à échéance aujourd'hui, à 16~\%
de \omterm{def:m1:volatility-first-contact:iv}{volatilité implicite}, tenue jusqu'à la clôture. La position rapporte si l'indice termine
à plus de 0.80~\% de son niveau de départ, ce qui, à cette volatilité, arrive
sur environ deux jours sur cinq. Données~: le tutoriel.}\label{fig:m1:zero-day-options-and-retail-flow:straddle}
\end{omfigure}

\section{Gamma près de l'échéance}

\begin{definition}[Couverture en delta et gamma]\label{def:m1:zero-day-options-and-retail-flow:gamma}
Le \emph{delta}\index{delta} d'une option est la variation de sa valeur pour une variation unitaire
du sous-jacent~; la \emph{couverture en delta}\index{couverture en delta} consiste à détenir la
quantité opposée du sous-jacent, de sorte que la position globale soit
insensible aux petits mouvements. Le \emph{gamma}\index{gamma} est la variation du delta
pour une variation unitaire du sous-jacent~: le rythme auquel la couverture doit être
ajustée. Une position longue d'options est longue en gamma~: son hedger vend à mesure que
le cours monte et achète à mesure qu'il baisse. Le hedger d'une position vendeuse d'options fait
l'inverse.
\end{definition}

\begin{proposition}[Le gamma à la monnaie croît comme $1/\sqrt{T}$]\label{prop:m1:zero-day-options-and-retail-flow:gamma}
Dans le modèle de Black--Scholes à taux nuls, le \omterm{def:m1:zero-day-options-and-retail-flow:gamma}{gamma} d'une option de
prix d'exercice $K$ vaut $\Gamma = \varphi(d_1)/(S\sigma\sqrt{T})$, avec $d_1 =
\ln(S/K)/(\sigma\sqrt T) + \sigma\sqrt T/2$ et $\varphi$ la densité normale. \omterm{def:m1:volatility-first-contact:atm}{À la
monnaie}, $\Gamma \approx 0.4/(S\sigma\sqrt{T})$~: diviser par deux le temps
restant jusqu'à l'échéance multiplie le \omterm{def:m1:zero-day-options-and-retail-flow:gamma}{gamma} par $\sqrt2$, et le \omterm{def:m1:zero-day-options-and-retail-flow:gamma}{gamma} est concentré à
environ un $\sigma\sqrt T$ du prix d'exercice.
\end{proposition}

\begin{omfigure}
\begin{tikzpicture}
\begin{axis}[width=10.5cm,height=4.6cm,xmode=log,x dir=reverse,
  xlabel={minutes de séance jusqu'à l'échéance (échelle logarithmique)},ylabel={points de delta pour 1~\%},
  tick label style={font=\small},label style={font=\small},
  xmin=4,xmax=10000,ymin=0,ymax=380,grid=major,grid style={gray!25},
  xtick={8190,1950,390,60,15,5},xticklabels={21~j,5~j,1~j,60,15,5},
  table/col sep=comma]
\addplot[omDef,very thick,mark=*,mark size=1.4pt] table[x=minutes,y=delta_change_pct]{figdata/markets-1/26-zero-day-options-and-retail-flow/atm_gamma.csv};
\draw[thin,dashed] (axis cs:4,100) -- (axis cs:10000,100);
\node[font=\footnotesize,anchor=south west] at (axis cs:9000,102) {100~: toute la position};
\end{axis}
\end{tikzpicture}

\omcaption{\omterm{def:m1:zero-day-options-and-retail-flow:gamma}{Gamma} \omterm{def:m1:volatility-first-contact:atm}{à la monnaie} à 16~\% de volatilité, exprimé comme variation du \omterm{def:m1:zero-day-options-and-retail-flow:gamma}{delta}, en
points de \omterm{def:m1:zero-day-options-and-retail-flow:gamma}{delta}, qu'une extrapolation linéaire donne pour un mouvement de 1~\%. Il franchit 100,
soit toute la position, à environ une heure de l'échéance~: au-delà, un mouvement de 1~\%
dépasse toute l'incertitude restante de l'option. Données~: la
construction du chapitre.}\label{fig:m1:zero-day-options-and-retail-flow:atm}
\end{omfigure}

\begin{omfigure}
\begin{tikzpicture}
\begin{axis}[width=10.5cm,height=4.8cm,
  xlabel={sous-jacent (prix d'exercice 100)},ylabel={variation du delta pour 1~\%},
  tick label style={font=\small},label style={font=\small},
  xmin=94,xmax=106,ymin=0,ymax=1.1,grid=major,grid style={gray!25},
  legend columns=3,legend style={font=\footnotesize,at={(0.5,-0.32)},anchor=north,draw=none,fill=none,/tikz/every even column/.append style={column sep=8pt}},
  table/col sep=comma]
\addplot[omExo,very thick,densely dotted] table[x=spot,y expr=\thisrow{month}/100]{figdata/markets-1/26-zero-day-options-and-retail-flow/gamma_profile.csv};
\addplot[omDef,very thick,dashed] table[x=spot,y expr=\thisrow{day}/100]{figdata/markets-1/26-zero-day-options-and-retail-flow/gamma_profile.csv};
\addplot[omThm,very thick] table[x=spot,y expr=\thisrow{hour}/100]{figdata/markets-1/26-zero-day-options-and-retail-flow/gamma_profile.csv};
\legend{un mois avant l'échéance,un jour,une heure}
\end{axis}
\end{tikzpicture}

\omcaption{\omterm{def:m1:zero-day-options-and-retail-flow:gamma}{Gamma} d'une option de prix d'exercice 100 à une volatilité de 16~\%, exprimé comme la
variation du \omterm{def:m1:zero-day-options-and-retail-flow:gamma}{delta} pour un mouvement de 1~\% du sous-jacent (extrapolation linéaire~: un
\omterm{def:m1:zero-day-options-and-retail-flow:gamma}{delta} ne peut pas varier de plus d'un). À un mois de l'échéance, la couverture varie de 0.09
pour 1~\%, quel que soit le cours~; à une heure de l'échéance, de la position entière
si le cours est au prix d'exercice et de rien s'il en est éloigné de 1~\%. Données~: la
construction du chapitre.}\label{fig:m1:zero-day-options-and-retail-flow:profile}
\end{omfigure}

À l'approche de l'échéance, le \omterm{def:m1:zero-day-options-and-retail-flow:gamma}{gamma} est tout ou rien. Le hedger d'une option dont le strike est égal
au cours actuel doit échanger la totalité de son notionnel, à l'achat puis à la vente, au fil
du franchissement du strike par le sous-jacent~; celui d'une option dont le strike est éloigné
de 1~\% conserve une position fixe et attend. Le \omterm{def:m1:option-contracts-and-exchanges:pin}{risque de pin} de la
\cref{def:m1:option-contracts-and-exchanges:pin} est la limite de ce tableau
à la clôture.

\section{Couverture des dealers et sous-jacent}

\begin{definition}[Gamma des dealers]\label{def:m1:zero-day-options-and-retail-flow:dealer}
Le \emph{gamma des dealers}\index{gamma des dealers} est le \omterm{def:m1:zero-day-options-and-retail-flow:gamma}{gamma} agrégé des
positions détenues par les \omterm{def:m1:options-market-structure:mm}{teneurs de marché d'options}, qui se couvrent, à la différence de leurs
clients, dont la plupart ne le font pas. Lorsque les dealers sont longs en \omterm{def:m1:zero-day-options-and-retail-flow:gamma}{gamma}, leur couverture
vend lors des hausses et achète lors des baisses, et tend à amortir les mouvements du sous-jacent~; lorsqu'ils sont courts en \omterm{def:m1:zero-day-options-and-retail-flow:gamma}{gamma}, elle
achète lors des hausses et vend lors des baisses, et tend à
les amplifier.
\end{definition}

\begin{omfigure}
\begin{tikzpicture}
\begin{axis}[width=10.5cm,height=4.8cm,
  xlabel={transactions des hedgers rapportées au mouvement précédent},ylabel={rapport de volatilité},
  tick label style={font=\small},label style={font=\small},
  xmin=-0.65,xmax=0.65,ymin=0.4,ymax=2.6,grid=major,grid style={gray!25},table/col sep=comma]
\addplot[omDef,very thick,mark=*,mark size=1.4pt] table[x=feedback,y=vol_ratio]{figdata/markets-1/26-zero-day-options-and-retail-flow/feedback.csv};
\draw[thin] (axis cs:0,0.4) -- (axis cs:0,2.6);
\node[font=\footnotesize,anchor=east] at (axis cs:-0.08,2.2) {hedgers longs en gamma};
\node[font=\footnotesize,anchor=west] at (axis cs:0.08,0.7) {hedgers courts en gamma};
\end{axis}
\end{tikzpicture}

\omcaption{Un marché jouet dans lequel chaque rendement de cinq minutes est une information plus une
fraction du rendement précédent, apportée par les transactions des hedgers~; l'axe vertical
est la \omterm{def:m1:volatility-first-contact:realised}{volatilité réalisée} rapportée à la volatilité de la seule information. Une rétroaction de
$-0.3$ retire un quart de la volatilité de la journée~; $+0.3$ en ajoute 40~\%. L'effet
est asymétrique~: l'amplification se cumule. Données~: la
simulation du tutoriel.}\label{fig:m1:zero-day-options-and-retail-flow:feedback}
\end{omfigure}

Le mécanisme ne fait pas débat~; sa taille et son signe un jour donné, si. Trois
précautions avant de croire à un chiffre d'«~exposition \omterm{def:m1:zero-day-options-and-retail-flow:gamma}{gamma}~». Premièrement, les
\omterm{def:m1:futures-contracts-and-exchanges:oi}{positions ouvertes} ne disent pas qui est long~: l'estimation usuelle \emph{suppose} que les
clients achètent des puts et vendent des calls, de sorte que les dealers sont courts des puts
et longs des calls, et son signe n'est que cette hypothèse. Deuxièmement, les options du
jour figurent à peine dans les \omterm{def:m1:futures-contracts-and-exchanges:oi}{positions ouvertes}, publiées une fois par jour à partir des
positions de la veille au soir~: les volumes ouverts et clos dans la séance ne laissent
aucune trace. Troisièmement, le flux des clients sur les options du jour est souvent
bilatéral, acheteurs et vendeurs des mêmes strikes se compensant en partie, de sorte que la
position \emph{nette} des dealers peut n'être qu'une faible fraction du volume. Un desk qui
veut ce chiffre doit l'estimer à partir des transactions signées, et non à partir des
\omterm{def:m1:futures-contracts-and-exchanges:oi}{positions ouvertes}.

\section{Flux d'options des particuliers}

Les investisseurs individuels représentent une part importante du volume d'options,
surtout sur les options à courte échéance portant sur quelques indices, fonds et grandes
actions. Leurs ordres atteignent le marché par l'intermédiaire des \omterm{def:m1:retail-flow-and-wholesaling:wholesaler}{wholesalers} et des
enchères présentées au \cref{ch:m1:options-market-structure}, et leurs caractéristiques sont
celles du flux actions des particuliers (\cref{ch:m1:retail-flow-and-wholesaling}), avec l'effet de
levier poussé à fond~: taille réduite, surtout des achats d'ouverture de calls et de puts hors
de la monnaie détenus pendant quelques heures ou quelques jours, concentration sur l'échéance
la plus proche, et peu d'information sur les minutes suivantes. Pour un \omterm{def:m1:what-a-trading-firm-does:market-maker}{teneur de marché}, ce
flux est précieux pour la même raison et risqué pour un nouvel entrant~: il est corrélé.
Lorsqu'une foule achète les mêmes calls sur la même action, les teneurs qui les ont vendus
sont courts en \omterm{def:m1:zero-day-options-and-retail-flow:gamma}{gamma} ensemble et achètent ensemble l'action à mesure qu'elle monte, dynamique
de la \cref{prop:m1:stock-loan-and-short-selling:cascade} avec des options à la place des
rachats
à découvert.

\section{Tutoriel~: gamma à la clôture}\label{tut:m1:zero-day-options-and-retail-flow:gamma}

\begin{tutorial}
\textbf{Objectif.} Calculer le \omterm{def:m1:zero-day-options-and-retail-flow:gamma}{delta} et le \omterm{def:m1:zero-day-options-and-retail-flow:gamma}{gamma} en temps de négociation, observer le \omterm{def:m1:zero-day-options-and-retail-flow:gamma}{gamma} \omterm{def:m1:volatility-first-contact:atm}{à la monnaie}
croître à mesure que l'échéance approche, traduire une position en actions à négocier
pour un mouvement de 1~\%, et simuler la rétroaction de la couverture sur le sous-jacent.
\textbf{État final~:} les quatre figures de données de ce chapitre.
\begin{tutsteps}
\item \textbf{\omterm{def:m1:zero-day-options-and-retail-flow:gamma}{Gamma}.} Une seule formule, la même pour les calls et les puts.
\omcode{code/firm/gex/firm_gex.py}{30}{36}{Gamma de Black--Scholes avec taux nuls.}\label{lst:m1:zero-day-options-and-retail-flow:gamma}
\item \textbf{Des contrats aux actions.} Linéaire, pour les petits mouvements~; et exact,
en réévaluant le \omterm{def:m1:zero-day-options-and-retail-flow:gamma}{delta}.
\omcode{code/firm/gex/firm_gex.py}{53}{61}{Actions pour un mouvement de 1~\%, et transaction de couverture exacte pour un mouvement donné.}\label{lst:m1:zero-day-options-and-retail-flow:shares}
\item \textbf{L'hypothèse de signe}, rendue explicite pour qu'on ne puisse pas
l'oublier.
\omcode{code/firm/gex/firm_gex.py}{64}{69}{Positions ouvertes vues du côté des dealers, sous une hypothèse énoncée.}\label{lst:m1:zero-day-options-and-retail-flow:dealer}
\item \textbf{Rétroaction.} Un rendement autorégressif dont le coefficient est la
transaction des hedgers par unité de mouvement.
\end{tutsteps}
\textbf{Pour aller plus loin.} Faire dépendre le coefficient de rétroaction de la distance entre le cours
et un strike à forte position ouverte, positif d'un côté d'un seuil et négatif de l'autre, et chercher le
pinning.
\end{tutorial}

\section{Construction~: l'estimateur d'exposition gamma}\label{bld:m1:zero-day-options-and-retail-flow:gex}

\begin{build}
\bfield{Rôle} Le \omterm{def:m1:options-market-structure:mm}{teneur de marché d'options} de la petite entreprise doit savoir, à tout moment,
combien d'actions son propre livre le forcera à négocier pour le prochain mouvement de 1~\%~;
son côté recherche veut la même estimation pour le marché, avec les
hypothèses étiquetées.
\bfield{Interface} \texttt{delta(spot, strike, years, vol, right)}; \texttt{gamma(spot, strike, years, vol)}; \texttt{Holding(strike, right,
years, vol, contracts, multiplier)}~; \texttt{net\_delta\_shares(book, spot)}; \texttt{gamma\_shares\_per\_pct(book, spot)}; \texttt{hedge\_trade(book,
spot\_before, spot\_after)}~; \texttt{dealer\_book\_from\_open\_interest(open\_interest, years, vol,
dealer\_side)}.
\bfield{Règles} Temps en minutes de négociation (252 jours de 390 minutes), car le temps
calendaire fausse le prix d'une option à deux heures de vie. Les options expirées ont un \omterm{def:m1:zero-day-options-and-retail-flow:gamma}{delta}
égal à un ou à zéro et pas de \omterm{def:m1:zero-day-options-and-retail-flow:gamma}{gamma}. Le nombre linéaire est accompagné de la transaction de
couverture exacte pour un mouvement énoncé à côté, car près de l'échéance les deux diffèrent
d'un facteur deux ou plus. L'estimation à l'échelle du marché prend sa convention de signe
comme argument explicite et n'a pas de valeur par défaut.
\bfield{Tests d'acceptation} \texttt{code/firm/gex/tests/}~: \omterm{def:m1:zero-day-options-and-retail-flow:gamma}{gamma} comme dérivée numérique du
\omterm{def:m1:zero-day-options-and-retail-flow:gamma}{delta}, égal pour les calls et les puts~; la loi en $1/\sqrt T$~; le sens de la couverture pour un
\omterm{def:m1:zero-day-options-and-retail-flow:gamma}{gamma} long et un \omterm{def:m1:zero-day-options-and-retail-flow:gamma}{gamma} court~; les options expirées~; deux conventions de signe donnant des
réponses opposées sur les mêmes \omterm{def:m1:futures-contracts-and-exchanges:oi}{positions ouvertes}.
\bfield{Pour aller plus loin} Classification des transactions signées à partir du flux d'options,
pour remplacer l'hypothèse, par codes de capacité client et \omterm{def:m1:what-a-trading-firm-does:market-maker}{teneur de marché} lorsque le flux
les fournit.
\end{build}

\begin{omsources}
\item Cboe, \emph{The State of the Options Industry: 2025}; communiqué de presse
de Cboe, ``Cboe to Add Tuesday and Thursday Expirations for SPX Weeklys
Options'', 13 avril 2022.
\end{omsources}

\section{Exercices}

\begin{exercise}[$\star$]\label{exo:m1:zero-day-options-and-retail-flow:1}
En utilisant $0.8\,\sigma\sqrt{T}$, calculer le prix du straddle \omterm{def:m1:volatility-first-contact:atm}{à la monnaie} arrivant
à échéance aujourd'hui sur un indice affichant une volatilité de 16~\%, exprimé en pourcentage de l'indice. Même question avec une volatilité de 32~\%.
\end{exercise}

\begin{exercise}[$\star$]\label{exo:m1:zero-day-options-and-retail-flow:2}
Un desk détient 500 calls \omterm{def:m1:volatility-first-contact:atm}{à la monnaie} (multiplicateur 100) sur une action cotée à 100 avec
un jour avant l'échéance et une volatilité de 16~\%. Le \omterm{def:m1:zero-day-options-and-retail-flow:gamma}{delta} du call vaut 0.50 et son \omterm{def:m1:zero-day-options-and-retail-flow:gamma}{gamma}
vaut 0.396 par dollar. Donner la couverture, ainsi que l'estimation linéaire du nombre d'actions à
négocier si l'action monte de 1~\%. Dans quel sens~?
\end{exercise}

\begin{exercise}[$\star$]\label{exo:m1:zero-day-options-and-retail-flow:3}
À partir de la \cref{prop:m1:zero-day-options-and-retail-flow:gamma}, de quel facteur
le \omterm{def:m1:zero-day-options-and-retail-flow:gamma}{gamma} \omterm{def:m1:volatility-first-contact:atm}{à la monnaie} croît-il entre un mois (21 jours) et un jour avant
l'échéance~? Entre un jour (390 minutes) et 15 minutes~?
\end{exercise}

\begin{exercise}[$\star\star$]\label{exo:m1:zero-day-options-and-retail-flow:4}
Pour la position de l'exercice~2, le \omterm{def:m1:zero-day-options-and-retail-flow:gamma}{delta} exact du call à 101 vaut 0.84.
Donner la transaction de couverture exacte pour la hausse de 1~\% et la comparer à
l'estimation linéaire. Pourquoi l'estimation linéaire surestime-t-elle la transaction~?
\end{exercise}

\begin{exercise}[$\star\star$]\label{exo:m1:zero-day-options-and-retail-flow:5}
Les options sur indice à échéance le jour même s'échangent à hauteur de 2.3 millions de contrats par jour, avec un multiplicateur de
100~\$ et un indice proche de 6\,000. Calculer le notionnel échangé. Un commentateur
conclut que «~1.4 trillion dollars of hedging hits the market every day~».
Donner deux raisons pour lesquelles cette affirmation est fausse de plusieurs ordres de grandeur.
\end{exercise}

\begin{exercise}[$\star\star$]\label{exo:m1:zero-day-options-and-retail-flow:6}
Le livre d'un \omterm{def:m1:options-market-structure:mm}{teneur de marché d'options} sur une action est vendeur en \omterm{def:m1:zero-day-options-and-retail-flow:gamma}{gamma} après une journée
d'achats de calls par la clientèle de détail. Décrire ses transactions de couverture si l'action
monte de 3~\% dans la dernière heure, puis baisse de 3~\% le lendemain matin. Qu'a-t-il gagné
ou perdu, qualitativement, par la seule couverture, et qu'a-t-il perçu en contrepartie~?
\end{exercise}

\begin{exercise}[$\star\star\star$]\label{exo:m1:zero-day-options-and-retail-flow:7}
\emph{Programmation.} Avec \texttt{simulate\_with\_hedgers(4000, 78, 0.001, f, 26)},
donner le ratio de volatilité pour $f = -0.3$, 0 et $+0.3$. Dériver analytiquement le ratio pour
un horizon long à partir de l'autorégression, et comparer.
\end{exercise}

\begin{exercise}[$\star\star\star$]\label{exo:m1:zero-day-options-and-retail-flow:8}
\emph{Trouver la faille.} Un service publie chaque matin~: «~Le \omterm{def:m1:zero-day-options-and-retail-flow:dealer}{gamma des dealers} vaut
$+\$4$ milliards par 1~\%~: la volatilité sera contenue aujourd'hui.~» Il calcule ce chiffre à partir
des \omterm{def:m1:futures-contracts-and-exchanges:oi}{positions ouvertes} de la veille, en supposant que les dealers sont acheteurs de tous les calls
et vendeurs de tous les puts. Citer trois raisons pour lesquelles ce chiffre peut avoir la mauvaise
taille ou le mauvais signe pour les options à échéance le jour même.
\end{exercise}

\section{Problème~: La dernière heure}

\begin{problem}[{Problème du week-end --- deux millions d'actions par point de pourcentage}]\label{pb:m1:zero-day-options-and-retail-flow:1}
Un dealer est vendeur de 20\,000 calls (multiplicateur 100) sur un \omterm{def:m1:exchange-traded-funds:etf}{fonds indiciel coté}
à 600, de prix d'exercice 600 et arrivant à échéance à la clôture d'aujourd'hui. La volatilité vaut 16~\%~; le temps est
mesuré en minutes de séance, 390 par jour et 252 jours par an. Le dealer
est couvert en \omterm{def:m1:zero-day-options-and-retail-flow:gamma}{delta}.

\medskip\noindent\textbf{Partie I --- La couverture maintenant.}
\begin{enumerate}
\item Donner le notionnel de la position en actions et en dollars.
\item Avec 60 minutes restantes, le \omterm{def:m1:zero-day-options-and-retail-flow:gamma}{delta} du call vaut environ 0.50. Donner la couverture.
\item Donner le mouvement d'un écart-type du fonds sur l'heure restante,
en pourcentage.
\item Le \omterm{def:m1:zero-day-options-and-retail-flow:gamma}{gamma} vaut 0.168 par dollar. Donner le nombre linéaire d'actions à négocier
par mouvement de 1~\%, avec son sens.
\item Comparer avec le notionnel. Que dit cette comparaison du
nombre linéaire~?
\end{enumerate}

\medskip\noindent\textbf{Partie II --- Transactions exactes.}
La construction donne les transactions de couverture exactes~: pour $+0.25\%$, acheter 472\,000 actions~; pour
$+1\%$, acheter 987\,000.
\begin{enumerate}[resume]
\item Donner le \omterm{def:m1:zero-day-options-and-retail-flow:gamma}{delta} du call après chaque mouvement.
\item Pourquoi la seconde transaction ne vaut-elle qu'environ le double de la première, pour un mouvement
quatre fois plus grand~?
\item À l'ouverture (390 minutes restantes), le nombre linéaire valait 792\,000 actions
par 1~\%. Avec 15 minutes restantes, il vaut 4.04 millions. Vérifier le rapport avec la valeur
à 60 minutes, à l'aide de la loi en $1/\sqrt T$.
\item Le fonds échange 60 millions d'actions par jour. Exprimer le nombre linéaire à 60 minutes
comme une fraction du volume moyen d'une heure.
\end{enumerate}

\medskip\noindent\textbf{Partie III --- La clôture.}
Le fonds est à 600.30 avec cinq minutes restantes.
\begin{enumerate}[resume]
\item Décrire la couverture du dealer au cours des cinq prochaines minutes si le cours oscille
entre 599.70 et 600.30.
\item Que coûte chaque oscillation au dealer~?
\item Qu'a perçu le dealer pour supporter ce risque, et de qui~?
\item Les options sont dénouées par livraison physique. À quoi d'autre le dealer doit-il
penser à 16:00~?
\item Si le dealer avait été \emph{acheteur} des 20\,000 calls, décrire ses transactions
et leur effet sur le cours du fonds vers 600.
\end{enumerate}

\medskip\noindent\textbf{Partie IV --- Jugement.}
\begin{enumerate}[resume]
\item Les options à échéance le jour même menacent-elles la stabilité des marchés~? Donner le mécanisme
par lequel elles pourraient le faire, et le fait relatif aux flux de la clientèle qui le limite.
\item Pourquoi les \omterm{def:m1:futures-contracts-and-exchanges:oi}{positions ouvertes} sont-elles un mauvais guide du positionnement intrajournalier~?
\item Un vendeur de straddles à échéance le jour même encaisse la prime trois jours sur
cinq. À quoi ressemble la distribution du P\&L, et qu'est-ce qui dimensionne la
position~?
\item Pourquoi les bourses apprécient-elles les échéances quotidiennes~?
\item Énoncer le \emph{résultat nommé}~: le nombre d'actions à négocier par mouvement de 1~\% avec
une heure restante, selon l'estimation linéaire.
\item En une phrase~: qu'est-ce qui distingue une option le jour de son échéance~?
\end{enumerate}
\end{problem}

\section{Questions d'entretien}

\begin{interviewq}[$\star$ \iqroles{trader, chercheur}]\label{iq:m1:zero-day-options-and-retail-flow:1}
Qu'est-ce que le \omterm{def:m1:zero-day-options-and-retail-flow:gamma}{gamma}, et que signifie être long ou short en \omterm{def:m1:zero-day-options-and-retail-flow:gamma}{gamma} pour
la façon dont vous couvrez~?
\end{interviewq}

\begin{interviewq}[$\star$ \iqroles{trader, chercheur}]\label{iq:m1:zero-day-options-and-retail-flow:2}
Comment le \omterm{def:m1:zero-day-options-and-retail-flow:gamma}{gamma} d'une option \omterm{def:m1:volatility-first-contact:atm}{à la monnaie} évolue-t-il à l'approche de l'échéance~?
Et celui d'une option hors de la monnaie~?
\end{interviewq}

\begin{interviewq}[$\star\star$ \iqroles{trader, chercheur}]\label{iq:m1:zero-day-options-and-retail-flow:3}
Évaluer de tête un straddle \omterm{def:m1:volatility-first-contact:atm}{à la monnaie} à échéance le jour même. Quel est son point mort~?
\end{interviewq}

\begin{interviewq}[$\star\star$ \iqroles{chercheur, ingénieur ML}]\label{iq:m1:zero-day-options-and-retail-flow:4}
Comment estimeriez-vous le positionnement en \omterm{def:m1:zero-day-options-and-retail-flow:dealer}{gamma des dealers}, et comment testeriez-vous
si votre estimation prédit quoi que ce soit~?
\end{interviewq}

\begin{interviewq}[$\star\star$ \iqroles{trader}]\label{iq:m1:zero-day-options-and-retail-flow:5}
Vous êtes short en \omterm{def:m1:zero-day-options-and-retail-flow:gamma}{gamma} à la clôture, le sous-jacent se trouvant à votre prix d'exercice.
Quels sont vos choix~?
\end{interviewq}

\begin{interviewq}[$\star\star\star$ \iqroles{chercheur, trader}]\label{iq:m1:zero-day-options-and-retail-flow:6}
Les options à échéance le jour même font-elles augmenter ou baisser la volatilité de l'indice~? Construisez
l'argument, puis dites quelles données permettraient de trancher.
\end{interviewq}
